\section*{Abstract}
		\noindent The T0 model presents an alternative theoretical framework for unifying fundamental physics. Starting from a single geometric constant $\xipar = \frac{4}{3} \times 10^{-4}$ and a universal energy field $\Efield(x,t)$, all physical phenomena are interpreted as manifestations of three-dimensional space geometry. The model replaces the 20+ free parameters of the Standard Model by the single parameter $\xipar$ (together with motivated integer settings) and offers deterministic explanations for quantum phenomena. Concrete predictions make the approach empirically testable. This treatise gives a non-technical overview of the theoretical foundations, mathematical structures, and experimental predictions.

	\medskip\noindent\textbf{Note on versions.} The German and English versions of this document differ in structure and extent; earlier revisions shortened, extended or changed parts of both. Where they differ, the more extensive German version applies; sections contained only in this English version count as supplements.


	\medskip\noindent\textbf{Status markers.} Statements marked \textbf{[S]} are \emph{speculative}: a hint or conjecture that has not yet been derived or numerically checked in the corpus.

	

	\section{Introduction: The Vision of Unified Physics}
	
	Imagine being able to explain all of physics -- from the smallest subatomic particles to the largest galaxy clusters -- with a single, simple idea. That's exactly what the T0 model attempts to achieve. Modern physics works with several theories that cannot easily be joined together; the T0 model proposes a simpler path.
	
	Today's physics resembles a house built by different architects: The ground floor (quantum mechanics) follows different rules than the first floor (relativity theory), and neither really fits with the attic (cosmology). Physicists must determine over twenty different numbers -- so-called free parameters -- from experiments, without knowing why these numbers have exactly these values. It's as if you needed twenty different keys to open all the doors in the house, without understanding why each lock is different.
	
	\begin{revolutionary}[Core Idea]
		The T0 model proposes: What if there were only one master key? A single number from which the others can be derived -- the geometric constant $\xipar = \frac{4}{3} \times 10^{-4}$. In the T0 model, this number is motivated by the geometry of the three-dimensional space in which we live.
	\end{revolutionary}
	
	The claim: This one number -- together with a few motivated integer settings, while SI anchors serve only for conversion into units -- should suffice to calculate many other numbers in physics: the lepton masses as steps of a ladder in $\xipar$ with fixed rational quantum numbers, converted into units at an anchor (Doc.~006), the strength of gravity, even the temperature of the universe \textbf{[S]}. It would be as if it turned out that all the seemingly random phone numbers in a phone book are built according to a single, hidden pattern.
	
	\section{The Geometric Constant $\xipar$: Foundation of the Model}
	
	\subsection{What is this number?}
	
	Imagine you're baking a cake. No matter how big the cake becomes, the ratio of ingredients stays the same -- for a good cake, you always need the right ratio of flour to sugar to butter. The geometric constant $\xipar$ is such a fundamental ratio for our universe.
	
	\begin{equation}
		\boxed{\xipar = \frac{4}{3} \times 10^{-4} = 0.0001333...}
	\end{equation}
	
	This number may seem small and unremarkable, but the T0 model does not regard it as random. The fraction 4/3 might be familiar from music -- it's the frequency ratio of a perfect fourth, one of the most harmonic intervals. More importantly: The factor 4/3 appears in the geometry of three-dimensional space.
	
	Think of a sphere -- the most perfect shape in space. Its volume is calculated with the formula $V = \frac{4}{3}\pi r^3$. Here, too, the factor 4/3 appears. The T0 model interprets this as meaning that this number is built into the structure of space.
	
	\subsection{Why is this number so important?}
	
	To understand why $\xipar$ is so fundamental, imagine the universe as a giant orchestra. In conventional physics, each instrument (each particle, each force) has its own, seemingly random tuning. Physicists must measure the tuning of each individual instrument without understanding why an electron has exactly this mass or why gravity is exactly this strong (or rather: this weak).
	
	\begin{important}
		The central thesis of the T0 model: All instruments in the universe's orchestra are tuned to a single pitch -- and this pitch is $\xipar$. 
		
		From this follows:
		\begin{itemize}
			\item The mass of an electron? A step of the ladder $r_i\,\xipar^{p_i}$ with fixed rational quantum numbers $r_i$, $p_i$, converted into units via an anchor (Doc.~006)
			\item The strength of gravity? Proportional to $\xipar^2$ (hence its weakness)
			\item The strength of the nuclear force? Proportional to $\xipar^{-1/3}$ (hence its strength)
		\end{itemize}
	\end{important}
	
	It would be as if it turned out that all seemingly different colors in the universe are just different mixtures of a single primary color.
	
	\section{The Universal Energy Field: The Only Fundamental Entity}
	
	\subsection{Everything is energy -- but differently than you think}
	
	Einstein taught us with his famous formula $E = mc^2$ that mass and energy are equivalent. The T0 model goes a step further and says: There is only energy. What we perceive as matter, as particles, as solid objects, are in this picture different vibration patterns of a single, all-permeating energy field.
	
	Imagine empty space not as nothing, but as a calm ocean. What we call ``particles'' are waves on this ocean. An electron is a small, very rapidly circling wave. A photon is a wave that runs across the ocean. A proton is a more complex wave pattern, like a whirlpool in water.
	
	\begin{equation}
		\boxed{\square \Efield = \left(\nabla^2 - \frac{1}{c^2}\frac{\partial^2}{\partial t^2}\right) \Efield = 0}
	\end{equation}
	
	This equation may look complicated, but it says something very simple: The energy field behaves like waves on a pond. It can oscillate, spread, interfere with itself -- and from all these behaviors emerges the apparent diversity of our world.
	
	\subsection{How does energy become an electron?}
	
	Think of a guitar string. When you pluck it, it doesn't vibrate arbitrarily, but in very specific patterns -- the overtones. Similarly, the universal energy field can't vibrate arbitrarily, but only in specific, stable patterns. We perceive these stable vibration patterns as particles:
	
	\begin{itemize}
		\item \textbf{An electron}: Imagine a tiny tornado of energy that constantly rotates around itself. This rotation is so stable that it can persist for billions of years.
		
		\item \textbf{A photon}: Like a wave on the sea that spreads in a straight line. Unlike the electron-tornado, this wave isn't trapped in one place but always moves at the speed of light.
		
		\item \textbf{A quark}: An even more complex pattern, like three intertwined vortices that stabilize each other.
	\end{itemize}
	
	The crucial point: There are no ``hard'' particles, no tiny billiard balls. Everything is motion, everything is vibration, everything is energy in different forms.
	
	\section{Quantum Mechanics Reinterpreted: Determinism Instead of Probability}
	
	\subsection{The end of randomness?}
	
	Quantum mechanics is considered the strangest theory in physics. It claims that nature is fundamentally random at the smallest scales -- that even God plays dice, as Einstein put it. A radioactive atom doesn't decay for a specific reason, but purely randomly. An electron isn't at a specific location, but ``smeared'' over many locations simultaneously until we measure it.
	
	The T0 model sees it differently: What we take for randomness is just our ignorance about the exact vibration patterns of the energy field. It's like rolling dice -- the throw appears random, but if you knew exactly the movement of the hand, air resistance, and all other factors, you could predict the result.
	
	\begin{tcolorbox}
		In the T0 model, the famous Schrödinger equation is no longer a probability calculation but describes how the real energy field evolves. The ``wave function'' isn't an abstract probability but the actual energy density of the field:
		\begin{equation}
			i\hbar \frac{\partial \Psi}{\partial t} = \hat{H}\Psi \quad \text{becomes} \quad i\hbar \frac{\partial \Efield}{\partial t} = \hat{H}_{\text{Field}}\Efield
		\end{equation}
	\end{tcolorbox}
	
	\subsection{The uncertainty relation -- newly understood}
	
	Heisenberg's famous uncertainty relation states that you can never know exactly both where a particle is and how fast it's moving. The more precisely you measure one, the more uncertain the other becomes. Physicists interpreted this as a fundamental limit of our knowledge.
	
	The T0 model sees it differently: Uncertainty isn't a knowledge limit but expresses that time and energy are two sides of the same coin:
	\begin{equation}
		\Delta E \cdot \Delta t \geq \frac{\hbar}{2}
	\end{equation}
	
	It's like with a musical note: To determine the pitch (frequency = energy) precisely, the tone must sound for a certain time. An ultra-short click has no defined pitch. That's not a measurement limitation, but a basic property of vibrations.
	
	\subsection{Schrödinger's cat lives -- and is dead}
	
	The most famous thought experiment in quantum mechanics is Schrödinger's cat: A cat in a box is simultaneously dead and alive until someone looks. That sounds absurd, and that's exactly what Schrödinger wanted to show.
	
	In the T0 model, the solution is simpler: The cat is never simultaneously dead and alive. The energy field is in a specific state, we just don't know it. If the field vibrates such that the radioactive atom has decayed, the cat is dead. If not, it lives. No mystery, no parallel worlds -- just our ignorance of the exact field vibrations.
	
	\subsection{Quantum entanglement -- the ``spooky'' phenomenon}
	
	Einstein called it ``spooky action at a distance'' -- quantum entanglement. When two particles are entangled, one knows immediately what happens to the other, no matter how far apart they are. Measure one particle as ``spin up'', the other is automatically ``spin down''. Immediately. Faster than light. This seems to violate everything we know about the maximum speed in the universe.
	
	The T0 model offers the following explanation \textbf{[S]}: The two particles are not separate at all. They're two bumps of the same wave in the energy field. Imagine a long rope that you hold in the middle and shake. Waves appear at both ends that are perfectly coordinated -- not because they communicate, but because they're part of the same vibration.
	
	\begin{equation}
		|\Psi_{\text{entangled}}\rangle = \frac{1}{\sqrt{2}}(|00\rangle + |11\rangle) \quad \Rightarrow \quad \Efield(x_1, x_2) = \Efield^{\text{coherent}}
	\end{equation}
	
	When you ``measure'' one bump (hold the rope at one point), that automatically determines what happens at the other end. No communication, no faster-than-light speed -- just the natural coherence of an extended wave.
	
	\subsection{Quantum computers -- why they work}
	
	Quantum computers are considered the future of computing technology. They use the strange properties of quantum mechanics -- superposition and entanglement -- to solve certain problems much faster than classical computers. But why do they work?
	
	\begin{experimental}
		In the T0 model, the answer is: A quantum computer directly manipulates the vibration patterns of the energy field. It uses the natural ability of the field to superpose many different vibration patterns simultaneously:
		
		\begin{itemize}
			\item \textbf{Deutsch algorithm}: Finds out with a single measurement whether a function is constant or balanced -- 100\% success even in the T0 model
			\item \textbf{Grover search}: Finds a needle in a haystack -- 99.999\% success rate in the deterministic T0 model
			\item \textbf{Shor factorization}: Breaks encryptions by finding periods -- works identically
		\end{itemize}
		
		The minimal deviations (0.001\%) are smaller than current practical measurement accuracy.
	\end{experimental}
	
	\section{The Unification of Quantum Mechanics, Quantum Field Theory and Relativity}
	
	\subsection{The great puzzle of modern physics}
	
	Modern physics has a problem -- actually several. We have three great theories, each of which works excellently on its own, but they don't fit together. It's as if we had three different maps of the same area that contradict each other at the edges.
	
	\textbf{Quantum mechanics} perfectly describes the world of atoms and molecules, but it completely ignores gravity. \textbf{Quantum field theory} extends quantum mechanics to high energies and can create and annihilate particles, but it produces infinite intermediate results that have to be handled by renormalization. And the \textbf{General Theory of Relativity} wonderfully explains gravity as curvature of spacetime, but it cannot easily be quantized -- so far nobody knows for certain how to properly describe quantum gravity.
	
	Since Einstein, physicists have been looking for a theory that unites all three. The T0 model proposes such a unification, with a simpler rather than a more complicated structure.
	
	\subsection{One field for everything}
	
	Instead of different fields for different particles (electron field, quark field, photon field, hypothetical graviton field), there's only one field in the T0 model -- the universal energy field. All seemingly different fields of quantum field theory are just different vibration modes of this one field:
	
	\begin{important}
		Imagine a concert hall. The different instruments (violin, trumpet, drums) produce different sounds, but they all vibrate in the same air. The air is the medium for all tones. Similarly, the universal energy field is the medium for all particles and forces:
		\begin{itemize}
			\item \textbf{Electromagnetism}: Transverse waves in the energy field (like light waves)
			\item \textbf{Weak nuclear force}: Local rotations of the energy field
			\item \textbf{Strong nuclear force}: Knots of the energy field that hold quarks together
			\item \textbf{Gravity}: The density of the energy field itself -- without additional particles
		\end{itemize}
	\end{important}
	
	\subsection{Gravity without gravitons}
	
	This is where it gets particularly interesting. Physicists have been searching for decades for ``gravitons'' -- hypothetical particles that transmit gravity, analogous to photons for electromagnetism. But nobody has ever found a graviton, and a perturbative quantization of gravity is not renormalizable.
	
	\begin{revolutionary}[Gravity in the T0 Picture]
		In the T0 picture, no gravitons are needed. Gravity isn't a force like the others, but a geometric effect of energy density:
		
		\begin{equation}
			\text{Spacetime curvature} = \frac{8\pi G}{c^4} \times \text{Energy density of the field}
		\end{equation}
		
		Where the energy field is denser, space curves more strongly. Mass is concentrated energy, so mass curves space. We perceive this curvature as gravity.
	\end{revolutionary}
	
	The gravitational constant $G$ is meant not to be an independent natural constant but to follow from the geometric constant: $G = \frac{\xipar^2}{4m_e}\,C_{\text{conv}}\,K_{\text{frak}}$ (Doc.~012); the factor $C_{\text{conv}}$ contains the SI conversion. The weakness of gravity is attributed to the fact that $\xipar^2$ is a tiny number. The gravitational coupling $Gm^2/(\hbar c)$ itself is $5.9\times10^{-39}$ (proton) or $1.8\times10^{-45}$ (electron); this value does not follow from $\xipar^2 = 1.8\times10^{-8}$ alone.
	
	\subsection{How do the open questions fit into the T0 picture?}
	
	In the T0 model, several major open questions of physics receive a common answer:
	
	\textbf{The hierarchy problem} -- Why is gravity so much weaker than the other forces? In the T0 model, the answer is simple: The strengths of all forces are powers of $\xipar$. The strong nuclear force has the strength $\xipar^{-1/3} \approx 20$, electromagnetism $\xipar^0 = 1$, the weak nuclear force $\xipar^{1/2} \approx 0.01$, and gravity $\xipar^2 \approx 0.00000001$. In this picture, the hierarchy is not fine-tuning but follows from geometry.
	
	\textbf{The infinities of quantum field theory} -- When physicists calculate the interaction of particles, they often get infinite values. These are handled by a procedure called renormalization. In the T0 model, these infinities are not expected to appear because the energy field has a natural minimal structure determined by $\xipar$.
	
	\textbf{The singularities} -- Black holes and the Big Bang lead to singularities in relativity theory -- points of infinite density where physics breaks down. In the T0 model, there are no real singularities. A black hole is simply a region of maximum energy field density, and the Big Bang would not occur in the static T0 picture \textbf{[S]}.
	
	\subsection{Quantum gravity in the T0 picture}
	
	The biggest unsolved problem of modern physics is quantum gravity. How does gravity behave at smallest scales? Nobody knows. Attempts to ``quantize'' gravity (turn it into a quantum theory) have so far not led to a generally accepted theory; approaches such as string theory work with additional spatial dimensions.
	
	\begin{important}
		The T0 model does not provide for a separate theory of quantum gravity \textbf{[S]}: gravity is meant to be already part of the quantized energy field. At small scales, the quantum fluctuations of the field dominate; at large scales, they average out to the smooth spacetime curvature we perceive as gravity.
		
		It's like with water: At the molecular level, you see individual H$_2$O molecules dancing around wildly (quantum level). At the macroscopic level, you see a smooth liquid (classical gravity). Both are the same phenomenon at different scales.
	\end{important}
	
	\section{Experimental Comparisons and Predictions}
	
	\subsection{The anomalous magnetic moment of the muon}
	
	The anomalous magnetic moment of the muon is one of the most precise measurements in physics.

	A muon is like a heavy electron -- it has the same properties but weighs 207 times more. When a muon circles in a magnetic field, it behaves like a tiny magnet. Physicists can measure the strength of this magnet to eleven decimal places and compare it with the Standard Model prediction.

	Until 2021, measured value and Standard Model prediction differed by $251(59) \times 10^{-11}$. With the Fermilab final result and the 2025 White Paper of the Muon~$g-2$ Theory Initiative, this deviation is no longer significant; it therefore serves neither as evidence nor as the target of a T0 contribution. The simple formula $\frac{\xipar}{2\pi}(m_\mu/m_e)^2$ does not supply such a contribution, because it gives $0.91$ and not $245\times10^{-11}$ (see Doc.~049). The current status and the assessment of the T0 statements on the muon $g-2$ are given in Doc.~382.

	
	\subsection{What we can still test}
	
	The T0 model makes many more predictions that can be tested in coming years:
	
	\textbf{Redshift newly understood}: Light from distant galaxies is redshifted -- its wavelength is stretched. The standard explanation: The universe is expanding. The T0 model says: Light loses energy traversing the energy field \textbf{[S]}. This energy loss is achromatic: it stretches all wavelengths by the same factor, as is also observed. There is therefore no wavelength dependence by which the two interpretations could be told apart; on these data they are indistinguishable (Doc.~267).
	
	\textbf{The tau lepton}: The heaviest of the three leptons (electron, muon, tau) is experimentally difficult to study. The T0 model predicts its anomalous magnetic moment: from the ratio of the $g-2$ differences of the three leptons, anchored to the measured values of electron and muon, it follows that $a_\tau \approx 1.281 \times 10^{-3}$, about $1.04\times10^{-4}$ above the Standard Model prediction (Doc.~382). A fixed $\xipar$ additive contribution following the pattern of the earlier muon explanation, by contrast, does not yield a viable tau value. Present measurement bounds on $a_\tau$ are still about 85 times wider than this difference; future experiments will test this.
	
	\textbf{Modified quantum entanglement}: In extremely precise Bell experiments, tiny deviations of 0.001\% from standard predictions should occur. That's at the limit of today's measurement technology, but not impossible.
	
	\subsection{Why these tests are important}
	
	Each of these predictions is a test of the entire T0 model. If even one of them is clearly wrong, the model must be revised or discarded. That's the strength of science -- theories must face reality.
	
	But if these predictions are confirmed? Then this would be a strong argument that a large part of physics can be derived from a single geometric constant -- a considerable simplification compared with the present state.
	
	\section{Cosmological Implications: An Eternal Universe}
	
	\subsection{No Big Bang -- no end}
	
	\textbf{[S]} The cosmological statements in this section are speculative. The present corpus deliberately leaves out the cosmic sector ($H_0$, $\Lambda$, dark energy), because the Standard Model does not derive it either (P39); this section should therefore be read as an outlook, not as a result.

	Standard cosmology describes a universe that has been expanding for about 13.8 billion years from a very hot, dense initial state -- the Big Bang -- and whose expansion continues.

	The T0 model considers a different picture: The universe would have no beginning and no end; it would be eternal and static. The observed redshift would then not be due to expansion but to the energy loss of light on its long journey through space.
	
	\begin{revolutionary}[Picture: Redshift through Energy Loss]
		Imagine standing at a foggy lake at night. The lights on the other shore appear fainter -- not because they're moving away from you, but because the light is altered on its way through the fog.

		In the T0 picture, the universe is similar: The ``fog'' is the omnipresent energy field. Light from distant galaxies loses energy (becomes redder), not because the galaxies are fleeing, but because its path through the $\xipar$ field is geometrically lengthened. This path lengthening stretches all wavelengths by the same factor; the redshift is achromatic:
		\begin{equation}
			1 + z = e^{\xipar x}
		\end{equation}
		The length scale of $x$ does not follow from $\xipar$ alone but is calibrated via the measured $H_0$ (Doc.~280).
	\end{revolutionary}


	\subsection{The cosmic microwave background -- explained differently}
	
	Everywhere in the universe, there's a weak microwave radiation with a temperature of 2.725 Kelvin -- the cosmic microwave background (CMB). The standard explanation: It's the cooled afterglow of the Big Bang.
	
	The T0 model says: It's the equilibrium temperature of the universal energy field. Every field has a natural temperature at which absorption and emission of energy are in equilibrium. For the $\xipar$ field, this is meant to be the measured 2.725 K.
	
	It's like the temperature in a cave deep underground -- the same everywhere, not because there was a Big Bang there, but because the system is in thermal equilibrium.
	
	\subsection{Dark matter and dark energy in the T0 picture}
	
	One of the greatest mysteries of modern cosmology: According to the standard model of cosmology, 95\% of the universe consists of dark matter and dark energy, which so far are inferred only from their gravitational effects. Galaxies rotate too fast (dark matter is needed to hold them together), and the universe is expanding at an accelerated rate (dark energy drives it apart).
	
	The T0 model attempts to do without both \textbf{[S]}:
	\begin{itemize}
		\item \textbf{Galaxy rotation}: The modified gravity through the energy field is meant to explain the rotation curves without additional matter
		\item \textbf{Accelerated expansion}: Is interpreted as a consequence of the redshift through energy loss, not as genuine acceleration
	\end{itemize}

	If this is confirmed, the assumption of additional, not directly observed components would no longer be needed.
	
	\subsection{A cyclic universe}
	
	If the universe is eternal, what happens with entropy? The second law of thermodynamics says that disorder always increases. After infinite time, the universe should end in heat death -- everything evenly distributed, no more structures.
	
	The T0 model proposes cycles for this \textbf{[S]}: Local regions of the universe go through phases of order and disorder, contraction and expansion, but globally everything remains in equilibrium. It's like an eternal ocean -- locally there are waves and whirlpools that arise and disappear, but the ocean as a whole persists.
	